\documentclass[11pt]{article}

\usepackage[utf8]{inputenc}
\usepackage[T1]{fontenc}
\usepackage{lmodern}
\usepackage{geometry}
\usepackage{amsmath,amssymb,amsfonts,amsthm,mathtools}
\usepackage{bm}
\usepackage{enumitem}
\usepackage{microtype}
\usepackage{hyperref}
\usepackage{titlesec}
\usepackage{booktabs}
\usepackage{array}
\usepackage{setspace}
\usepackage{mathrsfs}
\usepackage{physics}

\geometry{margin=1in}
\hypersetup{colorlinks=true, linkcolor=black, urlcolor=blue, citecolor=black}
\setlist[enumerate]{topsep=0.4em,itemsep=0.35em}
\setlist[itemize]{topsep=0.3em,itemsep=0.25em}
\titleformat{\section}{\large\bfseries}{\thesection.}{0.5em}{}
\titleformat{\subsection}{\normalsize\bfseries}{\thesubsection.}{0.5em}{}
\setstretch{1.06}

\newcommand{\Aether}{\AE{}ther}
\newcommand{\Aflow}{\AE{}ther-flow}
\newcommand{\TheoryName}{\Aflow{} Interpretation of Relativity}
\newcommand{\TheoryProgram}{\Aether{} / \Aflow{} framework}

\newcommand{\CompletedMetric}{g^{\sharp}}
\newcommand{\MatterSpace}{\Psi}

\newcommand{\Car}{\mathrm{car}}
\newcommand{\Cur}{\mathrm{cur}}
\newcommand{\Hom}{\mathrm{hom}}

\newcommand{\ForSpace}[1]{\mathcal I_{#1}^{\mathrm{for}}}
\newcommand{\PrimShell}{\mathcal S_{\mathrm{prim}}}
\newcommand{\MetricOut}{\mathscr G_{\mathrm{out}}}
\newcommand{\MatterOut}{\mathscr T_{\mathrm{out}}}

\newcommand{\SubSectionPackage}{\mathfrak V^{\mathrm{subsec}}}
\newcommand{\MicroSectionPackage}{\mathfrak W^{\mathrm{microsec}}}
\newcommand{\NanoSectionPackage}{\mathfrak X^{\mathrm{nanosec}}}
\newcommand{\PicoSectionPackage}{\mathfrak Y^{\mathrm{picosec}}}
\newcommand{\PicoMinSectionCore}{\mathfrak Y^{\mathrm{pic,sec}}}
\newcommand{\FemtoSectionPackage}{\mathfrak Z^{\mathrm{femtosec}}}
\newcommand{\AttoSectionPackage}{\mathfrak A^{\mathrm{attosec}}}
\newcommand{\ZeptoSectionPackage}{\mathfrak B^{\mathrm{zeptosec}}}
\newcommand{\ZeptoImageCore}{\mathfrak B^{\mathrm{zep,img}}}
\newcommand{\ZeptoMinSectionCore}{\mathfrak B^{\mathrm{zep,sec}}}

\newcommand{\ZeptoSectionSpace}[1]{\mathcal U_{#1}^{\mathrm{zeptosec}}}
\newcommand{\ZeptoSectionBody}[1]{\mathcal B_{#1}^{\mathrm{zeptosec}}}
\newcommand{\ZeptoSectionFunctional}[1]{\mathscr U_{#1}^{\mathrm{zeptosec}}}
\newcommand{\ZeptoSectionShell}[1]{\mathcal Z_{#1}^{\mathrm{zeptosec}}}

\newcommand{\YoctoImageCore}{\mathfrak C^{\mathrm{yoc,img}}}
\newcommand{\YoctoMinSectionCore}{\mathfrak C^{\mathrm{yoc,sec}}}

\newcommand{\YoctoSelectorPackage}{\mathfrak C^{\mathrm{yoc,sel}}}
\newcommand{\YoctoSelectorSpace}[1]{\mathcal D_{#1}^{\mathrm{yoc,sel}}}
\newcommand{\YoctoSelectorBody}[1]{\mathcal B_{#1}^{\mathrm{yoc,sel}}}
\newcommand{\YoctoSelectorFunctional}[1]{\mathscr L_{#1}^{\mathrm{yoc,sel}}}
\newcommand{\YoctoSelectorShell}[1]{\mathcal Z_{#1}^{\mathrm{yoc,sel}}}

\newcommand{\PreSelectorPackage}{\mathfrak D^{\mathrm{presel}}}
\newcommand{\PreSelectorSpace}[1]{\mathcal E_{#1}^{\mathrm{presel}}}
\newcommand{\PreSelectorBody}[1]{\mathcal B_{#1}^{\mathrm{presel}}}
\newcommand{\PreSelectorProj}[1]{\Pi_{#1}^{\mathrm{presel}}}
\newcommand{\PreSelectorFiber}[2]{\mathcal F_{#1}^{\mathrm{presel}}(#2)}
\newcommand{\PreSelectorFunctional}[1]{\mathscr P_{#1}^{\mathrm{presel}}}
\newcommand{\PreSelectorSelectorMin}[1]{\mathfrak p_{#1}^{\mathrm{presel,sel}}}
\newcommand{\PreSelectorShell}[1]{\mathcal Z_{#1}^{\mathrm{presel}}}

\newcommand{\PreSelectorMinSectionCore}{\mathfrak D^{\mathrm{presel,sec}}}
\newcommand{\PreSelectorMinSection}[1]{\mathfrak d_{#1}^{\mathrm{presel,sec}}}
\newcommand{\PreSelectorImgBody}[1]{\mathcal B_{#1}^{\mathrm{presel,img}}}
\newcommand{\PreSelectorImgProj}[1]{\Pi_{#1}^{\mathrm{presel,img}}}
\newcommand{\PreSelectorImgFunctional}[1]{\mathscr I_{#1}^{\mathrm{presel,img}}}
\newcommand{\PreSelectorImgShell}[1]{\mathcal Z_{#1}^{\mathrm{presel,img}}}

\newtheorem{definition}{Definition}
\newtheorem{proposition}{Proposition}
\newtheorem{remark}{Remark}
\newtheorem{corollary}{Corollary}
\newtheorem{theorem}{Theorem}

\title{\textbf{The \TheoryName{}}\\[0.4em]
\large Primitive-Reservoir Observer-Reduction\\
\large Selector-Generating Substrate Pre-Selector Dynamics\\
\large Collapse to an Observer-Relevant Pre-Selector Minimizer-Section Core}
\author{}
\date{}

\begin{document}

\maketitle

\begin{abstract}
The current primitive-reservoir observer-reduction frontier already moved the live burden below the current yocto-section-minimizer-section-generating substrate selector dynamics package \(\YoctoSelectorPackage\): the active theorem derived that selector package canonically from a deeper selector-generating substrate pre-selector dynamics package \(\PreSelectorPackage\). Under the recorded safeguard against recursive depth drift, however, the next honest move at that level is not another same-output deeper-origin relay that merely preserves the same pre-selector package. The sharper benchmark-facing question is whether the full recorded pre-selector package still carries benchmark-neutral presentational freedom once the induced selector package is treated as fixed.

The present note proves that it does. Relative to the already fixed selector body, selector functional, and exact selector shell, the only independent pre-selector datum still carried by the current package is the canonical minimizer section over the recorded selector body. The minimizer-image body is exactly the image of that section, the restricted projection is its inverse, the exact-shell image is its shell restriction, and the restricted pre-selector functional on that image is canonically induced from the recorded selector functional along the same section. Separate listing of those image-level constituents therefore does not add independent benchmark-facing freedom.

The benchmark boundary remains fixed throughout. The one operative metric is still exactly \(\CompletedMetric\), the ordinary matter route is unchanged, there is no second operative metric, the low-energy matter law is not enlarged, and no stronger promotion language is licensed. The positive result is therefore a genuinely smaller theorem-level collapse strictly inside the current pre-selector package: the live burden now contracts to an observer-relevant pre-selector minimizer-section core \(\PreSelectorMinSectionCore\). What remains open is to derive or justify that sharper core from still more primitive substrate structure, or else prove a still smaller collapse strictly inside it.
\end{abstract}

\section{Introduction}

The active derivational program of \TheoryName{} again reaches a stage where depth alone is not enough. The current active-primary theorem already showed that the full selector-generating substrate pre-selector dynamics package \(\PreSelectorPackage\) is not the primitive stopping point at present scope, because the live selector package \(\YoctoSelectorPackage\) is recovered there canonically from it. After that step, another deeper-origin relay that merely preserves the same pre-selector package is no longer the honest default move.

That theorem also exposes the next honest question. Inside the pre-selector package itself, what is genuinely independent observer-relevant data, and what is merely presentational bookkeeping once the recorded selector body, selector functional, and exact selector shell are already fixed? The present note takes the sharper internal-collapse route rather than another same-output deeper relay. It shows that the pre-selector image body, the restricted projection, the shell image, and the restricted pre-selector functional on that image are not independent pieces at current scope. They are all recovered from a single canonical minimizer section over the recorded selector body.

\paragraph{Framework and claim boundary.}
\TheoryProgram{} denotes the broader \Aether{} / \Aflow{} framework built on the claims that the \Aether{} is the underlying four-dimensional substrate of reality, the \Aflow{} is its intrinsic ordered motion, observed three-dimensional space is the local experiential slice of that deeper substrate, S-time is the experienced order of change arising from matter, light, and the \Aflow{}, and observed expansion is the three-dimensional appearance of deeper four-dimensional ordered motion. Within that framework, \TheoryName{} denotes the exact-closure theory in which the effective gravitational dynamics are adopted to be exactly Einsteinian with universal matter coupling. In the active sequence, \emph{adoption} means use of that established relativistic dynamics without claiming substrate derivation, while \emph{derivation} is reserved for a first-principles recovery from explicit substrate variables.

The present note stays strictly inside that boundary. It does not introduce a second operative metric, it does not enlarge the low-energy matter law, and it does not reopen the already-spent yocto-section minimizer-section-core, yocto-section minimizer-image-core, or full yocto-section presentations as if they were again the live burden. Its narrower benchmark-facing role is to isolate the only remaining independent datum still carried by the current pre-selector package itself.

\section{Fixed Inputs}

\begin{definition}[Recorded primitive continuation data]
\label{def:fixed}
Fix the following continuation data from the active primitive-reservoir observer-reduction line.
\begin{enumerate}
    \item For each sector label \(\sigma\in\{\Car,\Cur,\Hom\}\), a primitive forcing shell
    \[
    \ForSpace{\sigma}.
    \]
    \item The product shell
    \[
    \PrimShell
    \equiv
    \ForSpace{\Car}\times \ForSpace{\Cur}\times \ForSpace{\Hom}.
    \]
    \item The recorded metric and ordinary-matter outputs
    \[
    \MetricOut:\PrimShell\to\{\text{observer metrics}\},
    \qquad
    \MatterOut:\PrimShell\times\MatterSpace\to\{\text{observer stress tensors}\},
    \]
    satisfying
    \begin{equation}
    \MetricOut(\mathbf w)=\CompletedMetric,
    \qquad
    \MatterOut(\mathbf w,\psi)
    =
    T_{\mu\nu}^{\mathrm{matter}}[\CompletedMetric,\psi]
    \label{eq:fixedoutputs}
    \end{equation}
    for every \(\mathbf w\in\PrimShell\) and every matter configuration \(\psi\in\MatterSpace\).
\end{enumerate}
\end{definition}

\begin{definition}[Recorded selector benchmark base]
\label{def:selectorbase}
Fix \(\sigma\in\{\Car,\Cur,\Hom\}\). The active line already fixes the following benchmark-facing selector data:
\begin{enumerate}
    \item a finite-dimensional Euclidean selector state space
    \[
    \YoctoSelectorSpace{\sigma};
    \]
    \item a compact convex selector body
    \[
    \YoctoSelectorBody{\sigma}\subset\YoctoSelectorSpace{\sigma}
    \]
    with nonempty interior;
    \item a nonnegative smooth selector functional
    \[
    \YoctoSelectorFunctional{\sigma}:\YoctoSelectorSpace{\sigma}\to[0,\infty);
    \]
    \item a closed embedded exact selector shell
    \[
    \YoctoSelectorShell{\sigma}\subset\YoctoSelectorBody{\sigma};
    \]
    \item benchmark faithfulness, meaning preservation of the benchmark outputs \eqref{eq:fixedoutputs}, no second operative metric, no enlarged low-energy matter law, and use of this selector base as already fixed upstream data on the active line.
\end{enumerate}
\end{definition}

\begin{definition}[Recorded downstream chain below the pre-selector frontier]
\label{def:downstream}
Fix \(\sigma\in\{\Car,\Cur,\Hom\}\). The active line already fixes the following downstream data and consequences:
\begin{enumerate}
    \item the current observer-relevant yocto-section minimizer-section core \(\YoctoMinSectionCore\), the current observer-relevant yocto-section minimizer-image core \(\YoctoImageCore\), the current observer-relevant zepto-section minimizer-section core \(\ZeptoMinSectionCore\), the current observer-relevant zepto-section minimizer-image core \(\ZeptoImageCore\), the current atto-section package \(\AttoSectionPackage\), the current femto-section package \(\FemtoSectionPackage\), the current pico-section package \(\PicoSectionPackage\), the current observer-relevant pico-section minimizer-section core \(\PicoMinSectionCore\), the current nano-section package \(\NanoSectionPackage\), the current micro-section package \(\MicroSectionPackage\), the current sub-section package \(\SubSectionPackage\), and the previously recorded shell-transport consequences used on the active line;
    \item the previously recorded fact that, once the current selector package \(\YoctoSelectorPackage\) is available, the current observer-relevant yocto-section minimizer-section core \(\YoctoMinSectionCore\), the current observer-relevant yocto-section minimizer-image core \(\YoctoImageCore\), and all of the downstream data listed in item (1) follow unchanged.
\end{enumerate}
\end{definition}

\begin{definition}[Recorded selector-generating substrate pre-selector dynamics]
\label{def:presel}
Fix \(\sigma\in\{\Car,\Cur,\Hom\}\). The recorded current pre-selector dynamics package \(\PreSelectorPackage\) for sector \(\sigma\) consists of the following data.
\begin{enumerate}
    \item A finite-dimensional Euclidean pre-selector state space \(\PreSelectorSpace{\sigma}\), a compact convex pre-selector body
    \[
    \PreSelectorBody{\sigma}\subset\PreSelectorSpace{\sigma}
    \]
    with nonempty interior, and an affine surjection
    \[
    \PreSelectorProj{\sigma}:\PreSelectorSpace{\sigma}\to\YoctoSelectorSpace{\sigma}
    \]
    satisfying
    \[
    \PreSelectorProj{\sigma}\bigl(\PreSelectorBody{\sigma}\bigr)=\YoctoSelectorBody{\sigma}.
    \]
    For each \(y\in\YoctoSelectorBody{\sigma}\), write
    \[
    \PreSelectorFiber{\sigma}{y}
    \equiv
    \bigl(\PreSelectorProj{\sigma}\bigr)^{-1}(y)\cap\PreSelectorBody{\sigma}.
    \]
    Each such fiber is required to be nonempty, compact, and convex.
    \item A nonnegative smooth pre-selector functional
    \[
    \PreSelectorFunctional{\sigma}:\PreSelectorSpace{\sigma}\to[0,\infty)
    \]
    such that, for every \(y\in\YoctoSelectorBody{\sigma}\), the restriction of \(\PreSelectorFunctional{\sigma}\) to \(\PreSelectorFiber{\sigma}{y}\) is strictly convex and attains a unique minimizer. The resulting minimizer depends smoothly on \(y\); write it as
    \[
    \PreSelectorSelectorMin{\sigma}:\YoctoSelectorBody{\sigma}\to\PreSelectorBody{\sigma}.
    \]
    These data satisfy
    \[
    \PreSelectorProj{\sigma}\circ\PreSelectorSelectorMin{\sigma}
    =
    \mathrm{id}_{\YoctoSelectorBody{\sigma}},
    \]
    and the effective minimum value recovers the recorded selector functional:
    \[
    \YoctoSelectorFunctional{\sigma}(y)
    =
    \PreSelectorFunctional{\sigma}\bigl(\PreSelectorSelectorMin{\sigma}(y)\bigr)
    =
    \min_{u\in\PreSelectorFiber{\sigma}{y}}
    \PreSelectorFunctional{\sigma}(u)
    \qquad
    \text{for every }
    y\in\YoctoSelectorBody{\sigma}.
    \]
    \item A closed embedded exact pre-selector shell
    \[
    \PreSelectorShell{\sigma}\subset\PreSelectorBody{\sigma}
    \]
    satisfying
    \[
    \PreSelectorShell{\sigma}
    =
    \left\{
    u\in\PreSelectorBody{\sigma}:
    \begin{array}{l}
    \PreSelectorProj{\sigma}(u)\in\YoctoSelectorShell{\sigma},\\[0.2em]
    \PreSelectorFunctional{\sigma}(u)
    =
    \displaystyle\min_{v\in\PreSelectorFiber{\sigma}{\PreSelectorProj{\sigma}(u)}}
    \PreSelectorFunctional{\sigma}(v)
    \end{array}
    \right\},
    \]
    and
    \[
    \PreSelectorProj{\sigma}\big|_{\PreSelectorShell{\sigma}}
    :
    \PreSelectorShell{\sigma}\xrightarrow{\cong}\YoctoSelectorShell{\sigma}
    \]
    is a diffeomorphism.
    \item Benchmark faithfulness, meaning preservation of the benchmark outputs \eqref{eq:fixedoutputs}, no second operative metric, no enlarged low-energy matter law, and presentation as a deeper substrate-side source for the current selector package rather than as a replacement benchmark theory.
\end{enumerate}
\end{definition}

\begin{remark}
Definitions \ref{def:selectorbase} and \ref{def:presel} fix the current frontier package. The previous active-primary theorem already showed that benchmark-facingly relevant selector data factor through the deeper pre-selector package. The present note asks the sharper next question: once the selector benchmark base is already fixed, what still survives as independent benchmark-facing pre-selector data?
\end{remark}

\section{Observer-Relevant Pre-Selector Minimizer-Section Core}

\begin{definition}[Observer-relevant pre-selector minimizer-section core]
\label{def:sectioncore}
Fix \(\sigma\in\{\Car,\Cur,\Hom\}\) and a benchmark-faithful pre-selector package in the sense of Definition \ref{def:presel}. The \emph{observer-relevant pre-selector minimizer-section core} \(\PreSelectorMinSectionCore\) for sector \(\sigma\) consists of:
\begin{enumerate}
    \item the canonical minimizer section
    \[
    \PreSelectorMinSection{\sigma}
    \equiv
    \PreSelectorSelectorMin{\sigma}
    :
    \YoctoSelectorBody{\sigma}\to\PreSelectorBody{\sigma};
    \]
    \item benchmark faithfulness in the sense of Definition \ref{def:presel}(4).
\end{enumerate}
The image body
\[
\PreSelectorImgBody{\sigma}
\equiv
\PreSelectorMinSection{\sigma}\bigl(\YoctoSelectorBody{\sigma}\bigr),
\]
the shell restriction
\[
\PreSelectorMinSection{\sigma}\big|_{\YoctoSelectorShell{\sigma}}
:
\YoctoSelectorShell{\sigma}\to\PreSelectorShell{\sigma},
\]
the inverse projection on that image, and the restricted pre-selector functional induced there are understood as derived data rather than separately listed core data.
\end{definition}

\begin{remark}
Definition \ref{def:sectioncore} is strictly smaller than Definition \ref{def:presel} at current scope. Once the recorded selector body, selector functional, and exact selector shell are fixed, the pre-selector image body, restricted projection, shell image, and restricted functional on that image are all recovered from the single canonical minimizer section.
\end{remark}

\section{Collapse of the Current Pre-Selector Package}

\begin{proposition}[The canonical minimizer section is well defined]
\label{prop:section}
For every benchmark-faithful pre-selector package, the map \(\PreSelectorMinSection{\sigma}\) of Definition \ref{def:sectioncore} is a smooth section of \(\PreSelectorProj{\sigma}\), so that
\[
\PreSelectorProj{\sigma}\circ\PreSelectorMinSection{\sigma}
=
\mathrm{id}_{\YoctoSelectorBody{\sigma}}.
\]
\end{proposition}

\begin{proof}
Definition \ref{def:presel}(1) makes each fiber \(\PreSelectorFiber{\sigma}{y}\) nonempty, compact, and convex. Definition \ref{def:presel}(2) states that \(\PreSelectorFunctional{\sigma}\) has a unique minimizer on each such fiber and that the minimizer depends smoothly on \(y\). Definition \ref{def:presel}(2) therefore gives a well-defined smooth map
\[
\PreSelectorMinSection{\sigma}:\YoctoSelectorBody{\sigma}\to\PreSelectorBody{\sigma}.
\]
Because every minimizer lies in the fiber over the point at which it is selected, we have
\[
\PreSelectorProj{\sigma}\bigl(\PreSelectorMinSection{\sigma}(y)\bigr)=y
\]
for every \(y\in\YoctoSelectorBody{\sigma}\), which is the displayed identity.
\end{proof}

\begin{proposition}[The canonical minimizer section reconstructs the image body and projection]
\label{prop:reconstructbody}
For every benchmark-faithful pre-selector package,
\[
\PreSelectorImgBody{\sigma}
=
\PreSelectorMinSection{\sigma}\bigl(\YoctoSelectorBody{\sigma}\bigr),
\]
and the restricted projection
\[
\PreSelectorImgProj{\sigma}
\equiv
\PreSelectorProj{\sigma}\big|_{\PreSelectorImgBody{\sigma}}
:
\PreSelectorImgBody{\sigma}\to\YoctoSelectorBody{\sigma}
\]
is a homeomorphism with inverse \(\PreSelectorMinSection{\sigma}\).
\end{proposition}

\begin{proof}
The displayed identity is Definition \ref{def:sectioncore}. By Proposition \ref{prop:section},
\[
\PreSelectorProj{\sigma}\circ\PreSelectorMinSection{\sigma}
=
\mathrm{id}_{\YoctoSelectorBody{\sigma}}.
\]
Restricting \(\PreSelectorProj{\sigma}\) to \(\PreSelectorImgBody{\sigma}\) yields
\[
\PreSelectorImgProj{\sigma}\circ\PreSelectorMinSection{\sigma}
=
\mathrm{id}_{\YoctoSelectorBody{\sigma}}.
\]
Conversely, every \(u\in\PreSelectorImgBody{\sigma}\) has the form
\[
u=\PreSelectorMinSection{\sigma}(y)
\qquad
\text{for some }
y\in\YoctoSelectorBody{\sigma}.
\]
Hence
\[
\PreSelectorMinSection{\sigma}\bigl(\PreSelectorImgProj{\sigma}(u)\bigr)
=
\PreSelectorMinSection{\sigma}\bigl(\PreSelectorProj{\sigma}(u)\bigr)
=
\PreSelectorMinSection{\sigma}\bigl(\PreSelectorProj{\sigma}(\PreSelectorMinSection{\sigma}(y))\bigr)
=
\PreSelectorMinSection{\sigma}(y)
=
u.
\]
Therefore \(\PreSelectorImgProj{\sigma}\) and \(\PreSelectorMinSection{\sigma}\) are mutually inverse continuous maps, so the restricted projection is a homeomorphism with inverse \(\PreSelectorMinSection{\sigma}\).
\end{proof}

\begin{proposition}[The canonical minimizer section reconstructs the restricted functional]
\label{prop:reconstructfunctional}
For every benchmark-faithful pre-selector package, the restricted pre-selector functional on the minimizer image is recovered from the canonical minimizer section and the fixed selector functional by
\[
\PreSelectorImgFunctional{\sigma}\bigl(\PreSelectorMinSection{\sigma}(y)\bigr)
=
\YoctoSelectorFunctional{\sigma}(y)
\qquad
\text{for every }y\in\YoctoSelectorBody{\sigma},
\]
where
\[
\PreSelectorImgFunctional{\sigma}
\equiv
\PreSelectorFunctional{\sigma}\big|_{\PreSelectorImgBody{\sigma}}.
\]
Equivalently,
\[
\PreSelectorImgFunctional{\sigma}(u)
=
\YoctoSelectorFunctional{\sigma}\bigl((\PreSelectorMinSection{\sigma})^{-1}(u)\bigr)
\qquad
\text{for every }u\in\PreSelectorImgBody{\sigma}.
\]
\end{proposition}

\begin{proof}
By definition of \(\PreSelectorImgFunctional{\sigma}\),
\[
\PreSelectorImgFunctional{\sigma}\bigl(\PreSelectorMinSection{\sigma}(y)\bigr)
=
\PreSelectorFunctional{\sigma}\bigl(\PreSelectorSelectorMin{\sigma}(y)\bigr).
\]
Definition \ref{def:presel}(2) then gives
\[
\PreSelectorFunctional{\sigma}\bigl(\PreSelectorSelectorMin{\sigma}(y)\bigr)
=
\YoctoSelectorFunctional{\sigma}(y)
\]
for every \(y\in\YoctoSelectorBody{\sigma}\). This proves the first displayed identity. The equivalent formula on \(\PreSelectorImgBody{\sigma}\) follows from Proposition \ref{prop:reconstructbody}, which shows that every \(u\in\PreSelectorImgBody{\sigma}\) has the unique form \(u=\PreSelectorMinSection{\sigma}(y)\).
\end{proof}

\begin{proposition}[The shell image is exactly the shell restriction of the canonical minimizer section]
\label{prop:reconstructshell}
For every benchmark-faithful pre-selector package, if
\[
\PreSelectorImgShell{\sigma}
\equiv
\PreSelectorMinSection{\sigma}\bigl(\YoctoSelectorShell{\sigma}\bigr),
\]
then
\[
\PreSelectorImgShell{\sigma}
\subset
\PreSelectorShell{\sigma},
\]
and the shell restriction
\[
\PreSelectorMinSection{\sigma}\big|_{\YoctoSelectorShell{\sigma}}
:
\YoctoSelectorShell{\sigma}\xrightarrow{\cong}\PreSelectorImgShell{\sigma}
\]
is a diffeomorphism with inverse
\[
\PreSelectorImgProj{\sigma}\big|_{\PreSelectorImgShell{\sigma}}
:
\PreSelectorImgShell{\sigma}\xrightarrow{\cong}\YoctoSelectorShell{\sigma}.
\]
\end{proposition}

\begin{proof}
By Definition \ref{def:presel}(3),
\[
\PreSelectorProj{\sigma}\big|_{\PreSelectorShell{\sigma}}
:
\PreSelectorShell{\sigma}\xrightarrow{\cong}\YoctoSelectorShell{\sigma}
\]
is a diffeomorphism, and
\[
\PreSelectorMinSection{\sigma}\bigl(\YoctoSelectorShell{\sigma}\bigr)
\subset
\PreSelectorShell{\sigma}.
\]
Restricting the shell diffeomorphism to \(\PreSelectorImgShell{\sigma}\) gives
\[
\PreSelectorImgProj{\sigma}\big|_{\PreSelectorImgShell{\sigma}}
\circ
\PreSelectorMinSection{\sigma}\big|_{\YoctoSelectorShell{\sigma}}
=
\mathrm{id}_{\YoctoSelectorShell{\sigma}}.
\]
Conversely, if \(u\in\PreSelectorImgShell{\sigma}\), then \(u=\PreSelectorMinSection{\sigma}(y)\) for some \(y\in\YoctoSelectorShell{\sigma}\), so
\[
\PreSelectorMinSection{\sigma}\bigl(\PreSelectorImgProj{\sigma}(u)\bigr)
=
u.
\]
Thus the shell restriction of the canonical minimizer section and the restricted shell projection are inverse diffeomorphisms.
\end{proof}

\begin{proposition}[The pre-selector image-level presentation is reconstructed from the section core]
\label{prop:reconstructimage}
For every benchmark-faithful pre-selector package, the pre-selector image-level presentation consisting of the image body \(\PreSelectorImgBody{\sigma}\), the restricted projection \(\PreSelectorImgProj{\sigma}\), the restricted functional \(\PreSelectorImgFunctional{\sigma}\), and the shell image \(\PreSelectorImgShell{\sigma}\) is reconstructed from the smaller minimizer-section core \(\PreSelectorMinSectionCore\) together with the fixed selector body \(\YoctoSelectorBody{\sigma}\), selector functional \(\YoctoSelectorFunctional{\sigma}\), and exact selector shell \(\YoctoSelectorShell{\sigma}\).
\end{proposition}

\begin{proof}
By Proposition \ref{prop:reconstructbody}, the image body is the image of the canonical minimizer section and the restricted projection is its inverse. By Proposition \ref{prop:reconstructfunctional}, the restricted pre-selector functional on that image is uniquely determined by the canonical minimizer section together with the fixed selector functional. By Proposition \ref{prop:reconstructshell}, the shell image is exactly the shell restriction of that same section. Benchmark faithfulness is already part of Definition \ref{def:sectioncore}(2). Therefore each listed image-level constituent is recovered from Definition \ref{def:sectioncore} plus the fixed selector data.
\end{proof}

\begin{proposition}[All current benchmark-facing uses factor through the section core]
\label{prop:downstream}
For every benchmark-faithful pre-selector package, all current benchmark-facing uses of the pre-selector package \(\PreSelectorPackage\) on the active line factor through the minimizer-section core \(\PreSelectorMinSectionCore\) together with the fixed selector data. In particular, the recovery of the current selector package \(\YoctoSelectorPackage\), the current observer-relevant yocto-section minimizer-section core \(\YoctoMinSectionCore\), the current observer-relevant yocto-section minimizer-image core \(\YoctoImageCore\), the current observer-relevant zepto-section minimizer-section core \(\ZeptoMinSectionCore\), the current observer-relevant zepto-section minimizer-image core \(\ZeptoImageCore\), the current atto-section package \(\AttoSectionPackage\), the current femto-section package \(\FemtoSectionPackage\), the current pico-section package \(\PicoSectionPackage\), the current observer-relevant pico-section minimizer-section core \(\PicoMinSectionCore\), the current nano-section package \(\NanoSectionPackage\), the current micro-section package \(\MicroSectionPackage\), the current sub-section package \(\SubSectionPackage\), and the previously recorded shell-transport consequences does not require the separate ambient pre-selector presentation once the section core is fixed.
\end{proposition}

\begin{proof}
Proposition \ref{prop:reconstructimage} reconstructs the entire image-level pre-selector presentation from the section core and the fixed selector body, functional, and shell. Definition \ref{def:presel}(2)--(3) records that the current selector package is recovered from precisely that minimizer image, its induced restricted functional, and its exact shell transport over the fixed selector body. Definition \ref{def:downstream}(2) then records that, once the current selector package is available, the current observer-relevant yocto-section minimizer-section core \(\YoctoMinSectionCore\), the current observer-relevant yocto-section minimizer-image core \(\YoctoImageCore\), and all downstream data listed in Definition \ref{def:downstream}(1) follow unchanged. Therefore every current benchmark-facing use of \(\PreSelectorPackage\) factors through \(\PreSelectorMinSectionCore\) together with the fixed selector data, and the separate ambient pre-selector presentation is not needed to recover the recorded downstream consequences.
\end{proof}

\begin{proposition}[Separate full pre-selector presentation is benchmark neutral at current scope]
\label{prop:neutral}
For every benchmark-faithful pre-selector package, the separate presentation of the ambient pre-selector state space \(\PreSelectorSpace{\sigma}\), the off-image body directions
\[
\PreSelectorBody{\sigma}\setminus\PreSelectorImgBody{\sigma},
\]
the full pre-selector shell outside the image shell, and the unrestricted pre-selector functional away from the minimizer image does not add independent benchmark-facing freedom beyond the minimizer-section core \(\PreSelectorMinSectionCore\).
\end{proposition}

\begin{proof}
By Proposition \ref{prop:reconstructimage}, the image body, restricted projection, restricted functional, and shell image are reconstructed from the canonical minimizer section together with the fixed selector data. By Proposition \ref{prop:downstream}, all current benchmark-facing uses of the pre-selector package already factor through that reconstructed data. Therefore the larger ambient pre-selector presentation does not supply additional benchmark-facing information beyond the smaller minimizer-section core.
\end{proof}

\begin{proposition}[The benchmark package remains fixed]
\label{prop:benchmark}
Every observer-relevant pre-selector minimizer-section core preserves the same benchmark one-metric package \eqref{eq:fixedoutputs}. In particular, the operative metric remains exactly \(\CompletedMetric\), the ordinary matter route is unchanged, there is no second operative metric, and the low-energy matter law is not enlarged.
\end{proposition}

\begin{proof}
This is part of benchmark faithfulness in Definition \ref{def:sectioncore}(2). The present note only removes presentational redundancy inside the current pre-selector package. It does not alter the benchmark exact-closure package.
\end{proof}

\section{Main Theorem}

\begin{theorem}[Selector-generating substrate pre-selector dynamics collapse to an observer-relevant pre-selector minimizer-section core]
\label{thm:main}
Fix the primitive continuation data of Definition \ref{def:fixed}, the recorded selector benchmark base of Definition \ref{def:selectorbase}, the recorded downstream chain of Definition \ref{def:downstream}, and a benchmark-faithful pre-selector package in the sense of Definition \ref{def:presel}. Then, for each sector \(\sigma\in\{\Car,\Cur,\Hom\}\), the live benchmark-facing burden inside the current pre-selector frontier collapses from the full selector-generating substrate pre-selector dynamics package \(\PreSelectorPackage\) to the smaller observer-relevant pre-selector minimizer-section core \(\PreSelectorMinSectionCore\). More precisely:
\begin{enumerate}
    \item the canonical minimizer section over the fixed selector body is well defined by Proposition \ref{prop:section};
    \item the image body and restricted projection are reconstructed from that section by Proposition \ref{prop:reconstructbody};
    \item the restricted pre-selector functional on the minimizer image is reconstructed from the same section together with the fixed selector functional by Proposition \ref{prop:reconstructfunctional};
    \item the exact shell image is exactly the shell restriction of that section by Proposition \ref{prop:reconstructshell};
    \item the full pre-selector image-level presentation is therefore recovered from the section core and the fixed selector data by Proposition \ref{prop:reconstructimage};
    \item all current benchmark-facing uses of the pre-selector package factor through the section core by Proposition \ref{prop:downstream};
    \item the separate full pre-selector presentation adds no independent benchmark-facing freedom beyond the section core by Proposition \ref{prop:neutral};
    \item the benchmark boundary remains fixed by Proposition \ref{prop:benchmark};
    \item consequently the remaining honest burden below the previous pre-selector frontier is no longer the full selector-generating substrate pre-selector dynamics package \(\PreSelectorPackage\) as such. What remains is to derive or justify the sharper minimizer-section core \(\PreSelectorMinSectionCore\) from still more primitive substrate structure, or else prove a still smaller theorem-level collapse strictly inside that section core.
\end{enumerate}
\end{theorem}

\begin{proof}
Item (1) is Proposition \ref{prop:section}. Item (2) is Proposition \ref{prop:reconstructbody}. Item (3) is Proposition \ref{prop:reconstructfunctional}. Item (4) is Proposition \ref{prop:reconstructshell}. Item (5) is Proposition \ref{prop:reconstructimage}. Item (6) is Proposition \ref{prop:downstream}. Item (7) is Proposition \ref{prop:neutral}. Item (8) is Proposition \ref{prop:benchmark}. Item (9) is the routing consequence of the previous items together with the active safeguard against counting another same-output deeper-origin relay as new front-line progress when the live benchmark-relevant datum is unchanged.
\end{proof}

\begin{corollary}[What must be done next]
\label{cor:next}
After Theorem \ref{thm:main}, the next honest benchmark-facing manuscript must either:
\begin{enumerate}
    \item derive or justify the observer-relevant pre-selector minimizer-section core \(\PreSelectorMinSectionCore\) from still more primitive substrate structure in a way that changes benchmark-facing status;
    \item identify a genuinely smaller theorem-level observer-relevant datum strictly inside \(\PreSelectorMinSectionCore\) and prove a sharper collapse to it;
    \item do either of the previous two while preserving the same benchmark one-metric package and the same claim boundary.
\end{enumerate}
Another same-output deeper-origin relay that preserves this section core and therefore merely reproduces the same current pre-selector package \(\PreSelectorPackage\), the same current selector package \(\YoctoSelectorPackage\), the same current observer-relevant yocto-section minimizer-section core \(\YoctoMinSectionCore\), the same current observer-relevant yocto-section minimizer-image core \(\YoctoImageCore\), the same current observer-relevant zepto-section minimizer-section core \(\ZeptoMinSectionCore\), the same current observer-relevant zepto-section minimizer-image core \(\ZeptoImageCore\), and the same downstream observer-relevant atto-section, femto-section, pico-section, nano-section, micro-section, and sub-section consequences, another re-expansion back to the larger pre-selector image-level presentation or full pre-selector presentation as if those recovered constituents were still independently live, another reopening of the already-spent yocto-section core layers as if they were again the active burden, or any move that introduces a second operative metric, enlarges the ordinary matter law, or uses stronger promotion language does not satisfy the present burden. If a quantum continuation is pursued, it matters benchmark-facingly only insofar as it derives this same section core, collapses it further, or closes a benchmark derivation gate in a genuinely new way.
\end{corollary}

\begin{proof}
Theorem \ref{thm:main} shows that the separate full pre-selector package is no longer the minimal benchmark-facing datum at this stage. Therefore a subsequent manuscript must improve on the smaller section core rather than merely restating the larger pre-selector presentation.
\end{proof}

\section{Conclusion}

The positive result of this note is that the current selector-generating substrate pre-selector dynamics are not the minimal live datum on the active derivational line. Once one keeps the selector benchmark base fixed, the only independent pre-selector datum still required at benchmark-facing scope is the canonical minimizer section over the recorded selector body. The minimizer-image body is its image, the restricted projection is its inverse, the exact-shell image is its shell restriction, and the restricted pre-selector functional on that image is transported from the recorded selector functional along the same section. The larger ambient pre-selector presentation therefore does not add independent benchmark-facing freedom.

This preserves the same benchmark one-metric package and the same claim boundary. The result does not provide a completed first-principles derivation of GR from substrate microphysics, and it does not license stronger promotion language. Its narrower benchmark-facing gain is that the remaining honest burden now contracts again: not to the full pre-selector package, but to the smaller observer-relevant pre-selector minimizer-section core that still must be derived or justified from deeper substrate structure, or collapsed further if an even sharper honest datum can be isolated.

\end{document}
